\clearpage
\SourceLeaf{145}
\begin{center}
\textsc{Gauss} an \textsc{Schumacher}. Göttingen, 3. Dezember 1831.
\end{center}

\begin{center}
. . . . . . . . . . .
\end{center}

Die von Ihnen erwähnte Stelle in \textsc{Zachs} A.\,G.\,E. ist mir wohlbekannt. Die dort erwähnte Anwendung der M. der kl. Q. betrifft einen früher in derselben Zeitschrift abgedruckten Auszug aus \textsc{Ulugh Beighs} Zeitgleichungstafel, die zu manchen ganz kuriosen Ergebnissen geführt hatte. Diese Ergebnisse hatte ich \textsc{Zach} mit der Bemerkung mitgeteilt, dass ich dabei eine mir eigene, seit Jahren gebrauchte Methode benutzt habe, um Größen, die mit zufälligen Fehlern behaftet sind, auf willkürfreie und konsequente Weise zu kombinieren, ohne ihm jedoch das Wesen der Methode selbst mitzuteilen. Ich glaube, Ihnen schon einmal geschrieben zu haben, dass ich auf keinen Fall diese Stelle hervorheben werde, in der die Methode zum ersten Mal öffentlich angedeutet ist, und auch nicht wünsche, dass einer meiner Freunde es mit meiner Zustimmung tut. Das hieße \textls{anerkennen, als bedürfe} meine Angabe (Theoria Motus Corporum Coelestium), dass ich seit 1794 diese Methode vielfach gebraucht habe, einer Rechtfertigung, und dazu werde ich mich nie bereitfinden. Als \textsc{Olbers} bezeugte\HistoricalNote{[*)]}{[*)}{In seiner Abhandlung: Über den veränderlichen Stern im Halse des Schwans; v.\,\textsc{Lindenau} und \textsc{Bohnenberger}, Zeitschrift für Astronomie. Band II. Seite 192. September-Oktober 1816.]}, dass ich ihm 1802\HistoricalNote{[**)]}{[**)}{Sollte heißen »1803«.]} die ganze Methode mitgeteilt habe, war dies zwar gut gemeint. Hätte er mich aber vorher gefragt, so hätte ich es aufs Entschiedenste missbilligt.

\begin{center}
. . . . . . . . . . . . . . .
\end{center}

\begin{center}
\textsc{Gauss} an \textsc{Olbers}. Braunschweig, 30. Juli 1806.
\end{center}

\begin{center}
. . . . . . . . . . .
\end{center}

Herr\ \textsc{von Zach} schreibt mir noch, dass Sie sich zur Besprechung von \textsc{Legendres} Werk über die Kometenbahnen erboten hätten. Mit Vergnügen werde ich Ihnen also das mir von Herrn\ \textsc{von Zach} zugeschickte Exemplar nach Bremen
